\providecommand{\dirfiles}[1]{\par\smallskip\noindent
  {\footnotesize\raggedright\emph{Files:} #1.\par}\smallskip}
\providecommand{\dirP}{\mathbb{P}}
\providecommand{\dirE}{\mathbb{E}}
\providecommand{\dirU}{\mathrm{U}}
\section{Imaginary quadratic fields of class-group exponent six and seven}\label{sec:cls}

\providecommand{\clsp}{\mathfrak p}
\providecommand{\clsq}{\mathfrak q}
\providecommand{\clsXz}{X_0}

The \emph{exponent} of a finite abelian group is the least positive integer that
annihilates every element. Throughout this section $D<0$ is a fundamental discriminant,
$x=\abs{D}$, $K=\Q(\sqrt D)$, and $\chi=\chi_D=(D/\cdot)$ is the primitive odd quadratic
character of conductor $x$. We put
\[
  \clsXz=3.1\cdot10^{20}.
\]
Elsenhans, Kl\"uners and Nicolae (EKN) enumerated all imaginary quadratic fields with
$x\le\clsXz$ whose class group has small exponent, and bounded $x$ for each exponent under
the Riemann hypothesis for $\zeta_K$ \cite{Elsenhans2020Imaginary}. For the exponents six
and seven their bounds, $2.5\cdot10^{25}$ and $3.9\cdot10^{30}$, lie above $\clsXz$, and
the two intervals in between were left open. We show that both intervals are empty. The
results are \cref{cls:thm:six} (432 fields and 1{,}236 orders of exponent six) and
\cref{cls:thm:seven} (33 fields and 40 orders of exponent seven). With
\cref{cls:prop:small} they complete the census for every exponent from one to eight.

Each proof has three parts. A split prime $p$ gives a norm equation
$a^2+xb^2=4p^e$, which forbids small split primes and bounds the coefficient $b$. An
explicit formula with a signed test function converts the location of the zeros of
$L(s,\chi)$ into a lower bound for the total contribution of split primes. A finite
computation shows that the norm equations of two split primes are incompatible whenever
their coefficients $b$ are equal (exponent six) or have a ratio of small height (exponent
seven).

In the two censuses the input theorem \cref{src:thm:dgrh} enters in exactly three places,
each time for $\zeta$ and one quadratic Dirichlet $L$-function, through
$\zeta_K(s)=\zeta(s)L(s,\chi)$:
it discharges the hypothesis of the EKN discriminant bounds; it makes the zero side of the
explicit formula of \cref{cls:prop:weil} nonnegative; and it discharges the hypothesis of
the order classification of Aiken, Ross, Shvydkoy and Xue \cite{Aiken2026Boundary}. The
norm lemmas, the finite certificates and the reduced-form computations are unconditional.
The exhaustive EKN enumeration through $\clsXz$ is used as a published computation.

\subsection{Published inputs and the exponents below six}\label{cls:sec:inputs}
\sources{grh_transfers; grh-exponent-six-and-transfers; exponent_seven_and_grh_transfers}

We use the following statements of \cite{Elsenhans2020Imaginary}.
\begin{enumerate}[label=\textup{(E\arabic*)}]
\item\label{cls:it:census} (Theorem~1 and Table~1.) The list of imaginary quadratic fields
of class-group exponent at most eight with $x\le\clsXz$ is known. It contains
$9,56,17,203,27,432,33,778$ fields of exact exponent $1,2,\dots,8$ respectively.
\item\label{cls:it:bound} (Theorem~3 and Table~2.) If $\zeta_K$ satisfies the Riemann
hypothesis and $\Cl(K)$ has exponent $e$, then $x\le2.3\cdot10^{20}$ for $e=5$,
$x\le X_1:=2.5\cdot10^{25}$ for $e=6$, and $x\le3.9\cdot10^{30}$ for $e=7$. The analytic
input is the bound of Bach and Sorenson \cite{Bach1996Explicit} for the least rational prime
that splits in $K$, combined with the norm inequality of \cref{cls:lem:norm}.
\item\label{cls:it:complete} (Theorem~2.) Under the same hypothesis the lists in
\ref{cls:it:census} are complete for the exponents $1,2,3,4,5$ and $8$. The exponents six
and seven are excluded from this assertion.
\end{enumerate}
Since $\zeta_K(s)=\zeta(s)L(s,\chi_D)$, the hypothesis of \ref{cls:it:bound} and
\ref{cls:it:complete} holds for every $K$ by \cref{src:thm:dgrh}.

\begin{proposition}[Exponents one to five and eight; \stG]\label{cls:prop:small}
The numbers of imaginary quadratic fields whose class group has exact exponent
$1$, $2$, $3$, $4$, $5$ and $8$ are $9$, $56$, $17$, $203$, $27$ and $778$ respectively.
\end{proposition}

\begin{proof}
This is \ref{cls:it:census} and \ref{cls:it:complete}, with the hypothesis supplied by
\cref{src:thm:dgrh}.
\end{proof}

\begin{remark}[Exponent five]\label{cls:rem:five}
For $e=5$ the bound in \ref{cls:it:bound} lies inside the enumerated range, so no interval
remains. Of the 27 fields, 25 have class group $\Z/5\Z$, and the class group is
$(\Z/5\Z)^2$ for $D=-12451$ and $D=-37363$. An independent enumeration of the $75{,}988$
negative fundamental discriminants with $x\le250{,}000$, through their $11{,}689{,}615$
reduced forms, finds 196 fields whose class number is a nontrivial power of five, of which
exactly these 27 have exponent five; see \ref{cls:it:c-five} in \cref{cls:sec:six}.
\end{remark}

For orders we use \cite[Theorem~1.3(5),(6) and Section~6]{Aiken2026Boundary}: under ERH,
the tables A.6 and A.7 of that paper, of negative order discriminants with proper ideal
class group of exponent six, respectively seven, are complete if the interval left open
by \ref{cls:it:census} and \ref{cls:it:bound} contains no fundamental discriminant of that
exponent. The construction reduces all analytic hypotheses to the classification of the
underlying quadratic fields, for which \cref{src:thm:dgrh} supplies them. The integer
$\clsXz$ itself is not a fundamental discriminant.

\subsection{Euler's list, integer boxes and K3 surfaces}\label{dir:sec:idoneal}
\sources{grh-transfer; grh_transfer_analysis_2026-11-25}

A positive integer $n$ is \emph{idoneal} if the class group of primitive positive
definite binary quadratic forms of discriminant $-4n$ has exponent dividing two;
equivalently, every such form has an ambiguous class. The discriminant $-4n$ can belong
to a nonmaximal order.

\begin{theorem}[Euler's list is complete; \stG]\label{dir:thm:idoneal}
There are exactly $65$ idoneal numbers:
\begin{align*}
&1,2,3,4,5,6,7,8,9,10,12,13,15,16,18,21,22,24,25,28,30,33,37,\\
&40,42,45,48,57,58,60,70,72,78,85,88,93,102,105,112,120,130,133,\\
&165,168,177,190,210,232,240,253,273,280,312,330,345,357,385,408,\\
&462,520,760,840,1320,1365,1848.
\end{align*}
\end{theorem}

\begin{proof}
Weinberger's theorem \cite{Weinberger1973Exponents} concerns fundamental discriminants
and uses a zero-free interval near $1$ for real quadratic Dirichlet characters, which
Proposition~\ref{dir:prop:positive} supplies. Kani's reduction from orders to fields
\cite[Theorems~6 and~22, Corollary~23]{Kani2011Idoneal}, together with the finite list,
gives the statement.
\end{proof}

The possible extra fundamental discriminant in Weinberger's unconditional statement can
give two orders and hence two further idoneal numbers, not one; Kani's Corollary~23
records the distinction. The $65$ fundamental discriminants of exponent dividing two and
the $65$ idoneal numbers are different lists. The same completeness was drawn earlier
from a weaker zero-free half-plane in \cite{OpenAI2026Quasi}. Two finite checks
accompany the theorem: an enumeration of reduced forms for $1\le n\le10\,000$, storing a
nonambiguous reduced form for each rejected $n$, and an integer sieve of the
representation problem below through $n=5\,000\,000$.
\dirfiles{\artifact{grh_transfer_analysis_2026-11-25/data/idoneal_certificates.json},
\artifact{grh-transfer/checks/idoneal.py},
\artifact{grh-transfer/checks/idoneal-output.txt}}

A box with integer side lengths $x,y,z$ has surface area $2n$ with $n=xy+yz+zx$.

\begin{corollary}[Integer boxes; \stG]\label{dir:cor:boxes}
A positive integer $n$ fails to be of the form $xy+yz+zx$ with integers $0<x<y<z$
exactly when $n$ is idoneal. Hence every even number greater than $3696$ is the
surface area of a box with three different integer sides, and $3696$ is not. If equal
sides are allowed, exactly $18$ positive even numbers are not surface areas, namely the
doubles of
\[
 1,2,4,6,10,18,22,30,42,58,70,78,102,130,190,210,330,462 ,
\]
so that $924$ is the last missing even area.
\end{corollary}

\begin{proof}
Kani's Remark~35(b), crediting Rains, identifies the idoneal numbers with the failures
for $0<x<y<z$ \cite{Kani2011Idoneal}; apply Theorem~\ref{dir:thm:idoneal}. The classification
with repetitions is Theorem~1.1 of Borwein and Choi \cite{BorweinChoi2000Representations},
conditional on the same hypothesis.
\end{proof}

For example, the box $8\times9\times23$ has area $926$.

\begin{corollary}[The K3--Enriques obstruction list; \stG]\label{dir:cor:k3}
The list of $577$ idoneal genera of positive definite integral lattices of
Brandhorst, Sonel and Veniani is complete, and so is their list of $550$ co-idoneal
lattices occurring in the criterion for a complex algebraic K3 surface to cover an
Enriques surface. Exactly $65$ of the co-idoneal lattices have rank ten.
\end{corollary}

\begin{proof}
An integral positive definite lattice has a symmetric integer-valued pairing; its genus
consists of the isometry classes locally isometric to it over every $\Zp$ and over
$\R$. A genus is \emph{idoneal} if each of its members has a vector of square one. In
rank two such a vector $u$ splits off a copy of $[1]$, since $v-(v,u)u$ is integral and
orthogonal to $u$ for every $v$, and the determinant $n$ fixes the remaining factor
$[n]$. Hence an idoneal genus of rank two consists of the single class $[1]\oplus[n]$,
and the characterization of idoneal numbers by genera identifies the admissible
determinants with Euler's list, which is complete by Theorem~\ref{dir:thm:idoneal}.
Brandhorst, Sonel and Veniani
\cite[Remark~1.2 and Addendum~1.6]{BrandhorstSonelVeniani2025Idoneal} enumerate the
other ranks from this input: the numbers of idoneal genera of rank $r=1,\dots,13$ are
\[
 1,65,110,122,107,76,47,24,13,6,4,1,1,
\]
with total $577$ and none for $r>13$. This is the enumeration asked for in Kani's
Problem~45 \cite{Kani2011Idoneal}.

For the geometric statement put $\Lambda^-=U\oplus U(2)\oplus E_8(-2)$, where $U$ is
the hyperbolic plane, $E_8$ is positive definite and parentheses scale the pairing. Let
$T$ be the orthogonal complement of the N\'eron--Severi lattice in $H^2(X,\Z)$ of a K3
surface $X$. Keum's criterion asks for a primitive embedding $T\hookrightarrow\Lambda^-$
whose orthogonal complement contains no vector of square $-2$. Among the lattices that
embed primitively, the \emph{co-idoneal} ones are those for which every such embedding
has a vector of square $-2$ in its complement. By
\cite[Proposition~3.15]{BrandhorstSonelVeniani2025Idoneal} each complement of a
co-idoneal $T$ is $L(-2)$ for $L$ in a uniquely determined idoneal genus of rank
$12-\rank T$, and their Algorithm~4 enumerates the candidates $T$ from the genera and
discriminant forms. Thus rank two on the arithmetic side supplies rank ten on the
geometric side. Their Addendum~1.11 and Table~1 give, for $\rank T=2,\dots,11$, the
counts
\[
 3,8,17,41,76,107,122,110,65,1,
\]
with sum $550$ and none in higher rank. Both finite enumerations are the published
computations of \cite{BrandhorstSonelVeniani2025Idoneal}; the counts concern isometry
classes of lattices, and one lattice can occur for many surfaces.
\end{proof}

\subsection{Representations \texorpdfstring{$n=xy+yz+zx$}{n = xy+yz+zx} with a small
prime \texorpdfstring{$x+y$}{x+y}}\label{dir:sec:xyz}
\sources{grh-transfer}

For $n\ge1$ consider the strictly ordered positive integral solutions of
\begin{equation}\label{dir:eq:xyz}
 n=xy+yz+zx,\qquad 0<x<y<z .
\end{equation}
By Corollary~\ref{dir:cor:boxes} a solution exists for every $n>1848$. The results of this
subsection are quantitative and use only the Riemann hypothesis for $\zeta$ and for
real Dirichlet characters. All constants are absolute unless indicated.

\begin{lemma}[From a split prime to three coordinates; \stU]\label{dir:lem:split}
Let $\chi(k)=\legendre{-4n}{k}$ be the real Kronecker character modulo $4n$; it is
nonprincipal, since $\chi(-1)=-1$. Let $p\le\sqrt n$ be a prime.
\begin{enumerate}[label=\textup{(\roman*)}]
\item If $p$ is odd, $p\nmid 4n$ and $\chi(p)=1$, let $b$ be the unique integer with
 $0<b<p/2$ and $b^2\equiv-n\pmod p$. Then
 \[
  x=b,\qquad y=p-b,\qquad z=\frac{n+b^2}{p}-b
 \]
 is a solution of \eqref{dir:eq:xyz} with $x+y=p$.
\item Conversely, if \eqref{dir:eq:xyz} has a solution with $x+y=p$, then
 $p\nmid2n$ and $\chi(p)=1$.
\end{enumerate}
Different primes give different solutions.
\end{lemma}

\begin{proof}
In (i), $(x+y)(x+z)=n+x^2$ by construction, which is \eqref{dir:eq:xyz}; and
$x<y<z$ because $b<p/2$ and $n+b^2>p^2$. The pair sum recovers $p$. In (ii) the same
product identity gives $x^2\equiv-n\pmod p$ with $0<x<p/2$; hence $p\nmid 2n$ and
$\chi(p)=1$.
\end{proof}

\begin{lemma}[A uniform smoothed prime sum; \stG]\label{dir:lem:smooth}
Fix a nonnegative smooth function $w$ supported in $(1,2)$ with $\int w=1$. For a
Dirichlet character $\eta$ modulo $q$ and $X\ge2$ put
$S(\eta)=\sum_k\Lambda(k)\eta(k)w(k/X)$. If $\eta$ is principal or real, then
\[
 S(\eta)=\mathbf 1_{\eta=\chi_0}\,X+O_w\bigl(\sqrt X\log(2q)\bigr),
\]
where $\chi_0$ is the principal character modulo $q$.
\end{lemma}

\begin{proof}
First use the primitive character $\eta^*$ of conductor $f\mid q$ inducing $\eta$,
treating $f=1$ as the zeta function. Its Mellin integral has integrand
$-(L'/L)(s,\eta^*)X^s\widehat w(s)$ with
$\widehat w(s)=\int_0^\infty w(u)u^{s-1}\,\dd u$. Shift the contour from $\Re s=2$ to
$\Re s=-1/2$. The pole of zeta gives $X\widehat w(1)=X$. The possible trivial zero at
$0$ costs $O_w(1)$. At each nontrivial zero $\rho$ the residue is
$-X^{\rho}\widehat w(\rho)$, counted with multiplicity. Mellin integration by parts
gives $\abs{\widehat w(\sigma+it)}\ll_w(1+\abs t)^{-3}$ uniformly for
$-1/2\le\sigma\le2$. The number of zeros with $j\le\abs{\Im\rho}<j+1$ is
$O(\log(2f(j+2)))$ by the zero-counting formula for Dirichlet $L$-functions
\cite{BennettMartinOBryantRechnitzer2021Counting} and the classical estimate for zeta.
By \cref{src:thm:dgrh} every $\rho$ has real part $1/2$, so the absolutely convergent
zero sum is $O_w(\sqrt X\log(2f))$. The functional equation and the Euler product at
$\Re(1-s)=3/2$ give $L'/L\ll\log(2f(\abs t+2))$ on the new line, which contributes
$O_w(X^{-1/2}\log(2f))$. The contour shift is taken along heights avoiding zero
ordinates, using the local partial fraction formula for $L'/L$ and the same zero count.
Returning to modulus $q$ changes the sum by $O_w(\log q)$, since for each $p\mid q$ at
most one power of $p$ lies in $(X,2X)$.
\end{proof}

\begin{theorem}[Many representations with a small prime pair sum; \stG]
\label{dir:thm:xyz}
There exist $C,c>0$ such that for every sufficiently large $n$ and every real $X$ with
\[
 C\log^2(4n)\le X\le\tfrac12\sqrt n
\]
there are at least $cX/\log X$ solutions of \eqref{dir:eq:xyz} for which $x+y$ is a
prime in $(X,2X)$. In particular there are $\gg\log^2n/\log\log n$ solutions with
$x,y\ll\log^2n$, and $\gg\sqrt n/\log n$ solutions altogether.
\end{theorem}

\begin{proof}
Put $q=4n$ and let $\chi$ be as in Lemma~\ref{dir:lem:split}. On integers coprime to $q$ the
function $(\chi_0+\chi)/2$ is the indicator of $\chi=1$, and it vanishes on the other
integers. By Lemma~\ref{dir:lem:smooth}, after deleting higher prime powers,
\[
 \sum_{\substack{p\ \mathrm{prime}\\ \chi(p)=1}}(\log p)\,w(p/X)
 =\frac X2+O_w\bigl(\sqrt X\log(2q)+\sqrt X\log^2(2X)\bigr).
\]
Indeed each exponent $k\ge2$ contributes at most $(2X)^{1/k}\log(2X)$, and there are
$O(\log(2X))$ exponents. For $X\ge C\log^2(4n)$ the first error is at most $X/8$ once
$C$ is large, and the prime-power error is $o(X)$ as the lower endpoint tends to
infinity. Each prime contributes at most $\norm{w}_\infty\log(2X)$, so there are
$\gg X/\log X$ primes $p\in(X,2X)$ with $\chi(p)=1$. They are odd, coprime to $4n$ and
at most $\sqrt n$, and Lemma~\ref{dir:lem:split} turns them into distinct solutions. Take
$X=C\log^2(4n)$ for the small coordinates and $X=\sqrt n/3$ for the total count.
\end{proof}

\begin{theorem}[Several integers at once; \stG]\label{dir:thm:xyzmany}
For each fixed $k\ge1$ and all sufficiently large $N$, given
$n_1,\dots,n_k\in[N,2N]$ there are $\gg_k\log^2N/\log\log N$ primes $p\ll_k\log^2N$
such that every $n_i$ has a solution of \eqref{dir:eq:xyz} with $x_i+y_i=p$.
\end{theorem}

\begin{proof}
Inflate the $k$ quadratic characters $\legendre{-4n_i}{\cdot}$ to the modulus
$Q=4\prod_in_i$ and let $H$ be the group they generate. The function
$\abs H^{-1}\sum_{\eta\in H}\eta$ is the indicator of simultaneous splitting among
integers coprime to $Q$. Each member of $H$ is real, so by Lemma~\ref{dir:lem:smooth} the
smoothed prime-power sum of this function has main term $X/\abs H$ and error
$O_w(\sqrt X\log(2Q))$. Since $\abs H\le2^k$ and $\log Q=O_k(\log N)$, the choice
$X=C_k\log^2N$ with $C_k$ large absorbs the error. Removing prime powers as before and
applying Lemma~\ref{dir:lem:split} to each $n_i$ gives the claim.
\end{proof}

\begin{proposition}[A logarithmic obstruction; \stU]\label{dir:prop:xyzlog}
There are infinitely many $n$ for which \eqref{dir:eq:xyz} has no solution with
$x+y\le(1-o(1))\log n$, even if composite pair sums are allowed.
\end{proposition}

\begin{proof}
For $Y\to\infty$ put $M=4\prod_{3\le p\le Y}p$, the product over primes. Prescribe
$n\equiv1\pmod4$, choose $-n$ to be a quadratic nonresidue at every odd prime $p\le Y$,
and take the representative in $[M,2M)$ given by the Chinese remainder theorem. Every
integer $3\le a\le Y$ has an odd prime divisor at most $Y$ or is divisible by $4$, so
$x^2\equiv-n\pmod a$ has no solution; this rules out $a=x+y$, because
$(x+y)(x+z)=n+x^2$. The prime number theorem gives $\log n=(1+o(1))Y$.
\end{proof}

\begin{proposition}[A stronger obstruction at prime inputs; \stG]\label{dir:prop:xyzprime}
There are infinitely many primes $n$ such that every solution of \eqref{dir:eq:xyz}
satisfies
\[
 x+y>\bigl(1/\log4-o(1)\bigr)\log n\,\log\log n .
\]
\end{proposition}

\begin{proof}
Impose the signs of the preceding proof modulo $M=4\prod_{3\le p\le Y}p$, this time
seeking a prime $n$. The sign modulo $4$ and the quadratic signs at the odd primes give
$r=\pi(Y)$ independent binary conditions on the units modulo $M$, by the Chinese
remainder theorem. The indicator of the prescribed set on these units is the product of
$r$ terms $(1+\eps_j\eta_j)/2$ with $\eps_j\in\set{1,-1}$ and $\eta_j$ the corresponding
local quadratic characters inflated to modulus $M$. Extend the indicator by zero off
the units, so that the constant term of its expansion is the principal character
divided by $2^r$. Independence makes every other character in the expansion
nonprincipal, and the absolute values of all coefficients sum to one. By
Lemma~\ref{dir:lem:smooth} and the removal of prime powers, for $X\ge2$ the smoothed sum of
the indicator over primes $p$ with weight $(\log p)w(p/X)$ is
\[
 \frac X{2^r}+O_w\bigl(\sqrt X\,[\log(2M)+\log^2(2X)]\bigr).
\]
Take $X=4^r\log^6(2M)$. As $Y\to\infty$, $\log M\sim Y$ and $r\sim Y/\log Y$ by the
prime number theorem. After division by $\sqrt X$ the main term is $\log^3(2M)$,
whereas the error is $O_w(Y^2)$. Hence there is a prime $n\in(X,2X)$ with all
prescribed signs. The congruence argument of Proposition~\ref{dir:prop:xyzlog} rules out
$x+y\le Y$, and $\log n=(\log4+o(1))Y/\log Y$ gives the bound.
\end{proof}

\begin{proposition}[When the least splitting prime is the minimum; \stG]
\label{dir:prop:xyzmin}
For all sufficiently large squarefree $n\equiv1,2\pmod4$, the minimum of $x+y$ over the
solutions of \eqref{dir:eq:xyz} equals the least odd prime that splits in
$\Q(\sqrt{-n})$.
\end{proposition}

\begin{proof}
Put $a=x+y$ and $b=x$ in any solution. Then $a\mid n+b^2$ and $0<b<a/2$. Since $-n$ is
$2$ or $3$ modulo $4$, we have $4\nmid a$. If every odd prime $p\mid a$ divided $n$, it
would also divide $b$, and squarefreeness would give $\vp(n+b^2)=1$; the odd part of
$a$ would then divide $b$, although it is at least $a/2$. Thus some odd prime $p\mid a$
is coprime to $n$, and $-n$ is a square modulo $p$. This splitting prime satisfies
$p\le a$. Conversely, by Theorem~\ref{dir:thm:xyz} the least odd splitting prime is at most
$\sqrt n$ for large $n$, so it supplies a solution with this pair sum by
Lemma~\ref{dir:lem:split}.
\end{proof}

The congruence restriction matters: for $n\equiv3\pmod4$ and $n>15$ the triple
$(1,3,(n-3)/4)$ is always a solution.

\begin{remark}[Finding a witness]\label{dir:rem:xyzalg}
The program \artifact{grh-transfer/checks/represent.py} tests $a=x+y$ in increasing
order. For each $a$ it tries $0<b<a/2$ with $b^2\equiv-n\pmod a$ and $n+b^2>a^2$, and
uses the formula of Lemma~\ref{dir:lem:split} with $p$ replaced by $a$. Every solution of
\eqref{dir:eq:xyz} satisfies $3a^2<4n$, so exhaustive failure proves that there is
none; for a positive answer the first witness minimizes $x+y$. By Theorem~\ref{dir:thm:xyz},
for large $n$ at most $O(\log^4n)$ pairs $(a,b)$ are tested, which is deterministic
polynomial time in the bit length of $n$. The program uses neither the unknown
threshold of the asymptotic estimate nor a factorization of $n$.
\end{remark}

\subsection{Norm equations at split primes}\label{cls:sec:norm}
\sources{grh_transfers; grh-exponent-six-and-transfers; exponent_seven_and_grh_transfers}

\begin{lemma}[Norm equation; \stU]\label{cls:lem:norm}
Let $p$ be a rational prime that splits in $K$, let $\clsp\mid p$, and suppose
$[\clsp]^j=1$ in $\Cl(K)$ for some $j\ge1$. Then there are integers $a\ge0$ and $b\ge1$
with
\begin{equation}\label{cls:eq:norm}
  a^2+xb^2=4p^j .
\end{equation}
In particular $x\le4p^j$.
\end{lemma}

\begin{proof}
Write $\clsp^j=(\alpha)$ with $\alpha=(a+b\sqrt D)/2\in\OO_K$. The valuations of $\alpha$
at $\clsp$ and at $\bar\clsp$ differ, so $\alpha\notin\Q$ and $b\ne0$. Taking norms gives
$4p^j=a^2+xb^2$, and one may replace $a,b$ by their absolute values.
\end{proof}

\begin{corollary}[Residue witnesses; \stU]\label{cls:cor:witness}
Suppose $\Cl(K)$ is annihilated by $e\ge1$, and let $r$ be a prime with $4r^e<x$.
\begin{enumerate}[label=\textup{(\roman*)}]
\item The prime $r$ does not split in $K$.
\item If $r$ is odd, $c\ge1$ is an integer and $M=xc^2$, then $-M$ is not a nonzero
square modulo $r$.
\end{enumerate}
\end{corollary}

\begin{proof}
Part (i) is \cref{cls:lem:norm} with $j=e$. For (ii), suppose $-M=Dc^2$ is a nonzero
square modulo $r$. Then $r\nmid cD$ and $(D/r)=1$, so $r$ splits in $K$, contrary to (i).
\end{proof}

In the applications $M$ is produced from two norm equations without knowledge of $x$ or
$c$, and (ii) is tested for the small primes $r$. No extraction of the squarefree part of
$M$ is needed.

\subsection{An explicit-formula inequality}\label{cls:sec:weil}
\sources{grh-exponent-six-and-transfers; exponent_seven_and_grh_transfers}

\begin{proposition}[Signed test; \stG]\label{cls:prop:weil}
Let $g$ be a compactly supported real step function with nonzero $L^2$ norm, and put
\[
 F(u)=\frac{\int_{\R}g(v)g(v+u)\,\dd v}{\int_{\R}g(v)^2\,\dd v},
 \qquad h(t)=\int_{\R}F(u)e^{itu}\,\dd u .
\]
Let $C$ be the largest absolute slope of the piecewise linear function $F$. Then
\begin{equation}\label{cls:eq:weil}
 \log x>2\sum_{n\ge2}\frac{\chi(n)\Lambda(n)}{\sqrt n}\,F(\log n)+2-\frac{14}{9}\,C .
\end{equation}
\end{proposition}

\begin{proof}
The function $F$ is even, continuous and piecewise linear, with $F(0)=1$ and $\abs{F}\le1$.
Its Fourier transform $h$ is the squared modulus of the Fourier transform of $g$ divided
by $\norm{g}_2^2$, so $h\ge0$ on the real line. By the explicit Fourier integral of a step
function and its reflected copy, $h$ extends to an entire function with
$O((1+\abs{t})^{-2})$ decay on every fixed horizontal strip. For $\psi=\Gamma'/\Gamma$ the
Guinand--Weil formula for the odd primitive character $\chi$ reads
\begin{equation}\label{cls:eq:gw}
 \sum_{\rho}h\bigl((\rho-\tfrac12)/i\bigr)=
 \log\frac{x}{\pi}+\frac{1}{2\pi}\int_{\R}h(t)\,\Re\psi\bigl(\tfrac34+\tfrac{it}2\bigr)\,\dd t
 -2\sum_{n\ge2}\frac{\chi(n)\Lambda(n)}{\sqrt n}F(\log n),
\end{equation}
the sum running over the nontrivial zeros of $L(s,\chi)$; see
\cite[Lemma~9]{Carneiro2024Fourier}, with the Fourier conventions converted. The odd parity
determines the argument $3/4$ of the digamma function. By \cref{src:thm:dgrh} every
$(\rho-\tfrac12)/i$ is real, so the left side is nonnegative. This is the only use of the
location of zeros.

By the integral formula for the digamma function and Fourier inversion, the archimedean
term, including its factor $1/(2\pi)$, equals
\[
 \psi(\tfrac34)+2\int_0^\infty\frac{e^{-3u/2}}{1-e^{-2u}}\bigl(1-F(u)\bigr)\,\dd u
 \le\psi(\tfrac34)+\frac{14}{9}\,C .
\]
Indeed the difference $\Re\psi(\tfrac34+\tfrac{it}2)-\psi(\tfrac34)$ is an integral of
$1-\cos(tu)$ against a positive density, and its nonnegativity against $h(t)$ permits the
interchange of integrals. For the inequality use $0\le1-F(u)\le Cu$ for $u\ge0$, expand the
denominator geometrically, and bound
$\sum_{j\ge0}(2j+\tfrac32)^{-2}\le\tfrac49+\tfrac13$ by the integral test. Moreover
$\log\pi-\psi(\tfrac34)>2$: one has $\log\pi>1.1$, Euler's constant satisfies
$\gamma>H_6-\log7>\tfrac12$, the first two terms of the series for the digamma difference
give $\psi(1)-\psi(\tfrac34)>\tfrac13+\tfrac1{14}$, and $\psi(1)=-\gamma$. Inserting these
bounds into \eqref{cls:eq:gw} proves \eqref{cls:eq:weil}.
\end{proof}

\begin{definition}[Twelve-step windows]\label{cls:def:window}
Fix $Y>1$ and a vector $V=(v_1,\dots,v_6)$ of positive heights. Let $\delta=\log(Y)/12$,
and let $g=g_{Y,V}$ take the values
\[
  v_1,\ v_2,\ \dots,\ v_6,\ -v_6,\ \dots,\ -v_2,\ -v_1
\]
on twelve successive intervals of length $\delta$ beginning at zero, and vanish elsewhere.
Writing $(w_1,\dots,w_{12})$ for this signed vector, extended by zero, the associated $F$
satisfies $F(j\delta)=\sum_iw_iw_{i+j}/\sum_iw_i^2$ for integers $j\ge0$, is linear between
consecutive knots $j\delta$, and vanishes for $\abs{u}\ge\log Y$. For a prime $p\le Y$ put
\[
 R_p=2\sum_{p^k\le Y}\frac{(-1)^k}{p^{k/2}}F(k\log p),\qquad
 L_p=-\frac{4\log p}{\sqrt p}\,F(\log p),\qquad W_Y(p)=\max\set{0,L_p},
\]
and define the \emph{inert baseline} and the \emph{ramification constant}
\[
 S(Y)=2-\frac{14}9\,C+\sum_{p\le Y}R_p\log p,\qquad
 B_Y=\max\Bigl(0,\max_{p\le Y}R_p\Bigr).
\]
All quantities depend on $V$. We use the two vectors
\[
 V_0=(21,14,10,6,4,1),\qquad V_1=(1000,933,686,449,251,81).
\]
\end{definition}

The opposite signs at the two ends of $g$ make $F$ negative near the end of its support.
Although $h$ is nonnegative, an inert prime, for which $\chi(p)=-1$, can then contribute
positively to the right side of \eqref{cls:eq:weil}. The quantity $S(Y)$ is that right
side when every prime up to $Y$ is inert, and $L_p$ is the amount lost when a prime with
$p^2>Y$ splits instead.

\begin{lemma}[Baseline and losses; \stG]\label{cls:lem:baseline}
Let $Y$, $V$ be as in \cref{cls:def:window}, and suppose that every prime $p\le Y$ that
splits in $K$ satisfies $p^2>Y$. Then
\[
 \log x> S(Y)-\sum_{\substack{p\mid D\\ p\le Y}}R_p\log p
        -\sum_{\substack{p\le Y\\ p\text{ split}}}L_p
 \ \ge\ S(Y)-B_Y\log x-\sum_{\substack{p\le Y\\ p\text{ split}}}W_Y(p).
\]
\end{lemma}

\begin{proof}
Since $F(\log n)=0$ for $n\ge Y$, only prime powers $p^k\le Y$ occur in
\eqref{cls:eq:weil}. An inert prime has $\chi(p^k)=(-1)^k$ and contributes $R_p\log p$.
A ramified prime contributes nothing. A split prime contributes only through $k=1$, namely
$2(\log p)F(\log p)/\sqrt p=R_p\log p-L_p$. This gives the first inequality. For the
second, $\sum_{p\mid D}\log p\le\log x$ and $R_p\le B_Y$.
\end{proof}

\begin{remark}[Exact arithmetic]\label{cls:rem:exact}
All numerical inequalities for $S(Y)$, $B_Y$, $C$ and $L_p$ quoted below are verified
with rational arithmetic. The knot interval containing $\log(p^k)$ is located by comparing
the integers $(p^k)^{12}$ and $Y^j$. Logarithms are enclosed by reducing the argument to
$1\le z\le2$ and summing finitely many terms of the positive series
$2\operatorname{artanh}\bigl((z-1)/(z+1)\bigr)$ with its geometric tail. Integer square
roots bound $1/\sqrt{p^k}$ on both sides. No zero ordinate and no class group enters these
checks, and no floating-point comparison decides an inequality. The programs and their
outputs are \ref{cls:it:c-zero} for exponent six and \ref{cls:it:c-explicit} for exponent
seven.
\end{remark}

\subsection{The exponent-six census}\label{cls:sec:six}
\sources{grh-exponent-six-and-transfers; grh_transfer_synthesis_2026-11-25; grh_transfers}

\begin{theorem}[Exponent-six census; \stG]\label{cls:thm:six}
Exactly $432$ imaginary quadratic fields have ideal class group of exponent six. They are
the fields in the exponent-six row of \textup{\cite[Table~1]{Elsenhans2020Imaginary}}, and
the largest absolute discriminant among them is $5\,761\,140$. Exactly $1{,}236$ imaginary
quadratic orders have proper ideal class group of exponent six, of which $804$ are
nonmaximal. The largest absolute order discriminant is $43\,522\,752$.
\end{theorem}

By \ref{cls:it:census} and \ref{cls:it:bound} it suffices, for the fields, to show that no
$K$ with $\clsXz<x\le X_1$ has class group annihilated by six. The finite part of the
argument is the following certificate.

\begin{lemma}[Pairs of sixth-power norm equations; \stU]\label{cls:lem:sixpairs}
Let $p<q\le5500$ be primes with $4p^6>\clsXz$, put $N=q^6-p^6$, and let $u$ be a positive
divisor of $N$ with
\[
   q^3-p^3<u\le\floor{\sqrt N}.
\]
If $M=4p^6-(N/u-u)^2$ exceeds $\clsXz$, then there is a prime $r$ with $19\le r\le151$
such that $-M$ is a nonzero square modulo $r$.
\end{lemma}

\begin{proof}
This is a finite computation. Factor
\[
 N=(q-p)(q+p)(q^2+pq+p^2)(q^2-pq+p^2).
\]
Each factor is at most $3\cdot5500^2$, so trial division suffices. Generate the divisors
from their prime exponents, test the two inequalities, and reduce $-M$ modulo the small
primes. Two independently written programs perform the search \ref{cls:it:c-pairs}: one
with arbitrary-precision integers and sets of nonzero squares, the other with 128-bit
products, a check of each factorization, and precomputed residue masks for the equivalent
residue $(v-u)^2-4p^6$, $v=N/u$. The counts, with multiplicity over prime pairs and
divisors, are as follows.
\begin{center}
\begin{tabular}{lr}\toprule
Quantity & Count\\\midrule
Primes $p\le5500$ with $4p^6>\clsXz$ & 414\\
Unordered pairs & 85{,}491\\
Divisors in the interval & 10{,}144{,}433\\
Pairs and divisors giving $M>\clsXz$ & 9{,}393{,}302\\
Candidates with no witness & 0\\\bottomrule
\end{tabular}
\end{center}
The second program includes the lower endpoint $u=q^3-p^3$, which adds one divisor with
$M=0$ for each pair and accounts for its count $10{,}229{,}924$ in the third row. Both
programs agree on the candidates above the threshold and report no survivor. A complete
replay of both is recorded in \ref{cls:it:c-replay}.
\end{proof}

\begin{proof}[Proof of \cref{cls:thm:six}]
Suppose $\clsXz<x\le X_1$ and $\Cl(K)$ is annihilated by six.

\emph{Split primes are large.} Since $4\cdot2000^6<\clsXz$, \cref{cls:cor:witness}(i)
shows that every prime that splits in $K$ exceeds $2000$.

\emph{There are at least $28$ split primes below $5500$.} Take $Y=5500$ and $V=V_0$ in
\cref{cls:def:window}, and let $N_Y$ be the number of split primes $p\le Y$. Each of them
satisfies $p^2>Y$, so \cref{cls:lem:baseline} applies. For any $L>0$ that bounds $L_p$
from above for all primes $2000<p\le5500$ it gives
\[
 N_YL>S(Y)-(1+B_Y)\log x\ge S(Y)-(1+B_Y)\log X_1 .
\]
The rational enclosures of \cref{cls:rem:exact} give the two strict inequalities
\begin{equation}\label{cls:eq:sixmargins}
 \max_{2000<p\le5500}L_p<0.216,\qquad S(Y)-(1+B_Y)\log X_1>5.84 .
\end{equation}
For orientation, $\sum_{p\le Y}R_p\log p$ is about $65.682$, $B_Y$ about $0.04476$ and
$C$ about $0.47355$. Since $27\cdot0.216=5.832<5.84$, we get $N_Y\ge28$.

\emph{There are at most $18$ coefficients.} By \cref{cls:lem:norm} every split prime
$p\le5500$ gives a norm equation
\begin{equation}\label{cls:eq:norm6}
   a_p^2+xb_p^2=4p^6,\qquad a_p\ge0,\quad b_p\ge1 .
\end{equation}
The exact inequality $4\cdot5500^6<19^2\clsXz$ gives $b_p\le18$. By the pigeonhole
principle two distinct split primes $p<q\le5500$ have the same coefficient $b=b_p=b_q$.

\emph{Eliminating the discriminant.} The integers $a_p<a_q$ have the same parity, since
their squares differ by a multiple of four. Subtract the two equations
\eqref{cls:eq:norm6} and put
\[
   u=(a_q-a_p)/2,\qquad v=(a_q+a_p)/2,\qquad N=q^6-p^6=uv .
\]
Then $0<u\le v$, $a_p=v-u$, and the positive integer
\[
   M=4p^6-(v-u)^2=xb^2
\]
is greater than $\clsXz$; in particular $4p^6>\clsXz$. Since $N=(q^3-p^3)(q^3+p^3)$ and
$N/u-u$ decreases for $u>0$, positivity of $M$ is equivalent, for $u\le\sqrt N$, to
$u>q^3-p^3$. Hence $(p,q,u)$ satisfies the hypotheses of \cref{cls:lem:sixpairs}, which
supplies a prime $r\le151$ with $-M$ a nonzero square modulo $r$. As
$4r^6<4\cdot2000^6<x$, this contradicts \cref{cls:cor:witness}(ii) with $c=b$.

Hence no such $K$ exists, and \ref{cls:it:census} gives the statement on fields. The
statement on orders now follows from \cite[Theorem~1.3(5)]{Aiken2026Boundary}, whose
remaining condition has just been discharged. Table~A.6 of that paper, in the form of the
authors' data file \cite{Ross2026Data}, has $1{,}236$ distinct entries with largest
absolute value $43\,522\,752$, and its subset of fundamental discriminants agrees entry
for entry with the list of $432$ fields; the other $1{,}236-432=804$ orders are nonmaximal.
\end{proof}

The $432$ field discriminants, with class number, number of ramified primes, largest prime
factor and cofactor, are listed in \ref{cls:it:c-fields}. The $1{,}236$ order discriminants
with their groups, a copy of \cite{Ross2026Data}, and the comparison of the two lists are
in \ref{cls:it:c-orders}. The distribution of the fields by the number $\omega(x)$ of
ramified primes is as follows.
\begin{center}
\begin{tabular}{lrrrrrrr}\toprule
$\omega(\abs{D})$ & 2&3&4&5&6&7&Total\\\midrule
Fields of exact exponent six&60&141&140&76&12&3&432\\\bottomrule
\end{tabular}
\end{center}
The three fields with seven ramified primes have $D=-3\,892\,980$, $-4\,696\,692$ and
$-5\,761\,140$. The largest value of $x/P^+(x)$ among the $432$ fields, where $P^+$
denotes the largest prime factor, is $204{,}204$.

\begin{remark}[Verification by reduced forms]\label{cls:rem:forms}
Membership in both lists has been checked independently of the cited tables. A negative
discriminant $\Delta$ of an order gives one proper ideal class for each reduced primitive
positive definite binary quadratic form $(a,b,c)$ with $b^2-4ac=\Delta$; the forms are
enumerated using $3a^2\le\abs{\Delta}$, and their number is the order of the class group.
Powers are computed as follows. Change variables unimodularly until $a$ is prime to
$\Delta$; such a representative exists by avoiding the bad residues at the primes dividing
$\Delta$. Let $\omega$ generate the order, with minimal polynomial $f$. The form
corresponds, up to the inverse convention, to the invertible ideal $(a,\omega-u)$ for a
simple root $u$ of $f$ modulo $a$. Lift $u$ to a root $u_j$ of $f$ modulo $a^j$ by
Hensel's rule. The ideal $(a^j,\omega-u_j)$ contains the $j$th power of $(a,\omega-u)$ and
has the same norm, hence equals it, and reduction of its norm form tests principality. A
finite abelian group annihilated by six whose order is divisible by six has exponent
exactly six, and its order determines its structure, because its Sylow subgroups are
elementary abelian. In this way the $432$ fields were recomputed through $x=5\,761\,140$
\ref{cls:it:c-fields}, and the listed group orders and the sixth power of each of the
$103\,830$ classes in the $1{,}236$ orders were verified \ref{cls:it:c-orders}.
Exhaustiveness rests on \ref{cls:it:census}, on \cref{cls:thm:six} and on the conductor
theorem of \cite{Aiken2026Boundary}.
\end{remark}

The delivered certificates for this subsection and for \cref{cls:rem:five} are the
following.
\begin{enumerate}[label=\textup{(C\arabic*)},series=clscert,leftmargin=*]
\item\label{cls:it:c-zero} Rational enclosures for \eqref{cls:eq:sixmargins}, program and
output:\\
\artifact{grh-exponent-six-and-transfers/checks/zero_test.py}\\
\artifact{grh-exponent-six-and-transfers/data/zero-test.txt}
\item\label{cls:it:c-pairs} The two searches of \cref{cls:lem:sixpairs} and their outputs:\\
\artifact{grh-exponent-six-and-transfers/checks/paired_norms.py}\\
\artifact{grh-exponent-six-and-transfers/data/paired-norms.txt}\\
\artifact{grh-exponent-six-and-transfers/checks/paired_norms.cpp}\\
\artifact{grh-exponent-six-and-transfers/data/paired-norms-cpp.txt}
\item\label{cls:it:c-replay} Replay of \ref{cls:it:c-zero} and \ref{cls:it:c-pairs}:\\
\artifact{grh_transfer_synthesis_2026-11-25/validation/full_replay/summary.json}
\item\label{cls:it:c-fields} The $432$ fields and the log of their recomputation:\\
\artifact{grh_transfers/checks/exponent_six.csv}\\
\artifact{grh_transfers/checks/exponent_six_run.txt}
\item\label{cls:it:c-orders} The $1{,}236$ orders, the comparison with the field list, and
the check of all class powers:\\
\artifact{grh-exponent-six-and-transfers/data/orders-six.txt}\\
\artifact{grh-exponent-six-and-transfers/data/order-check.txt}\\
\artifact{grh-exponent-six-and-transfers/data/order-class-check.txt}
\item\label{cls:it:c-five} The exponent-five enumeration of \cref{cls:rem:five}:\\
\artifact{grh_transfers/checks/exponent_five.csv}\\
\artifact{grh_transfers/checks/exponent_five_run.txt}
\end{enumerate}

\subsection{The exponent-seven census}\label{cls:sec:seven}
\sources{exponent_seven_and_grh_transfers; grh-exponent-six-and-transfers}

\begin{theorem}[Exponent-seven census; \stG]\label{cls:thm:seven}
Exactly $33$ imaginary quadratic fields have ideal class group of exponent seven. For
$31$ of them the group is cyclic of order seven. For the discriminants $-63499$ and
$-118843$ it is $(\Z/7\Z)^2$. Exactly $40$ imaginary quadratic orders have proper ideal
class group of exponent seven, of which seven are nonmaximal. The discriminants are those
of \cref{cls:tab:seven}.
\end{theorem}

At exponent seven the pigeonhole argument of \cref{cls:thm:six} is not available: with
the same test the forced number of split primes does not exceed the number of admissible
coefficients. Two further ingredients replace it. The elimination is carried out for
pairs of coefficients with a prescribed reduced ratio, and the upper bound for the total
loss charges each group of mutually excluded coefficients only once.

\subsubsection*{Arithmetic restrictions on split primes}
Let $D=-x$ be a fundamental discriminant with $x>\clsXz$ and suppose $\Cl(K)$ has
exponent seven. This group is a vector space over $\F_7$, so its order is odd. Genus theory
gives $2$-rank $t-1$ for $t$ distinct prime discriminant factors, so there is just one
ramified prime. The dyadic fundamental discriminants with this property are small, hence
$x$ is an odd prime with $x\equiv3\pmod4$. In particular every prime $p\le7000$ is
unramified. By \cref{cls:lem:norm} every split prime $p$ gives
\begin{equation}\label{cls:eq:norm7}
 a_p^2+xb_p^2=4p^7,\qquad a_p\ge0,\quad b_p\ge1 ,
\end{equation}
and by \cref{cls:cor:witness}(i) every prime $r$ with $4r^7\le\clsXz$, that is every
prime $r\le691$, is inert. This includes $2$, so $D\equiv5\pmod8$ and $x\equiv3\pmod8$.

\begin{lemma}[Finite exclusions; \stU]\label{cls:lem:sevenpairs}
Let $K$ be as above and let $p\ne q$ be primes that split in $K$.
\begin{enumerate}[label=\textup{(\roman*)}]
\item If $p,q\le6500$, then $b_p\ne b_q$.
\item If $p,q\le4200$, then the fraction $b_p/b_q$ in lowest terms does not have both
numerator and denominator at most $12$.
\end{enumerate}
\end{lemma}

\begin{proof}
(i) Take $p<q$ and suppose $b_p=b_q$. By \eqref{cls:eq:norm7}, $a_p$ and $a_q$ have the
same parity. Put $u=(a_q-a_p)/2$ and $v=(a_q+a_p)/2$. Then $uv=q^7-p^7$, $0<u\le v$, and
$M=4p^7-(v-u)^2=xb_p^2>\clsXz$. For each pair of primes $693<p<q\le6500$ the program
\ref{cls:it:c-equal} factors $q-p$ and $(q^7-p^7)/(q-p)$, enumerates all divisors $u\le\sqrt{q^7-p^7}$ of their product, and for
every divisor with $M>\clsXz$ searches for an odd prime $r$ with $4r^7<\clsXz$ such that
$-M$ is a nonzero square modulo $r$. By \cref{cls:cor:witness}(ii) such an $r$ excludes
the candidate. There are $717$ eligible primes, $256{,}686$ pairs and $191{,}603$ divisor
candidates above the threshold, and every candidate has a witness $r\le107$.

(ii) Suppose the contrary. By (i) the two coefficients differ, so, naming the primes
suitably, $b_p=ic$ and $b_q=jc$ with coprime integers $1\le i<j\le12$ and $c\ge1$. No
order between $p$ and $q$ is assumed, and both assignments of the two primes occur in the
enumeration. Put
\[
   d=i^2q^7-j^2p^7 .
\]
It is nonzero, since $p,q>693$. Multiplying the two equations \eqref{cls:eq:norm7} by
$j^2$ and $i^2$ and subtracting shows that the nonnegative integers $ia_q$ and $ja_p$
have squares differing by $4d$, so their sum and difference are even. The half sum $v$
and the half absolute difference $u$ satisfy $uv=\abs{d}$. If $d>0$ then $ja_p=v-u$ and
$ia_q=v+u$; if $d<0$ the two expressions are exchanged. For each divisor $u\le\sqrt{\abs d}$
of $\abs d$ one tests divisibility by $j$ and $i$, recovers the two integers $a_p,a_q$, and
keeps only the common integral values
\[
        M=\frac{4p^7-a_p^2}{i^2}=\frac{4q^7-a_q^2}{j^2}>\clsXz .
\]
Here $M=xc^2$, so \cref{cls:cor:witness}(ii) applies again: a prime $r>12$ with
$4r^7<\clsXz$ for which $-M$ is a nonzero square excludes the candidate. The search is
organized as follows.

The program \ref{cls:it:c-ratio} first discards the cases with $4p^7\le\clsXz i^2$ or $4q^7\le\clsXz j^2$, which contradict
\eqref{cls:eq:norm7}. It then tests the two norm equations simultaneously modulo
\[
             128,\ 81,\ 125,\ 49,\ 121,\ 13,\ 17,\ 19,\ 23 .
\]
For an odd base prime $r$ of one of these moduli choose a unit $w$ with $(-w/r)=-1$. Since
$r$ is inert, $(-x/r)=-1$, so $x/w$ is a square modulo $r$, and a square root lifts to each
$r^k$ by Hensel's lemma. At $2$ use $w=3$: a dyadic unit is a square precisely when it is
$1\bmod8$, and $x\equiv3\pmod8$. Thus every possible $M=xc^2$ lies in the finite set
$w(\Z/r^k\Z)^2$ of square multiples of $w$, including multiples of $r$. For each modulus a
pair $(p,q)$ with ratio $i/j$ is kept if there is a residue $t$ in this set such that
$4p^7-i^2t$ and $4q^7-j^2t$ are both squares; this condition is necessary, and the
residues $t$ may be chosen separately at the different moduli. For each surviving pair
every divisor $u\le\sqrt{\abs d}$ of $\abs d$ is tested as described. The results are as
follows.
\begin{center}
\begin{tabular}{lrr}\toprule
 & Equal coefficients & Distinct coefficients\\\midrule
Prime pairs processed after preliminary tests & 256{,}686 & 1{,}963{,}144\\
Divisor candidates giving a possible $M>\clsXz$ & 191{,}603 & 231{,}471\\
Largest first witness $r$ & 107 & 97\\
Candidates left & 0 & 0\\\bottomrule
\end{tabular}
\end{center}
The first column counts unordered pairs. The second counts ordered prime pairs for each
of the $45$ reduced ratios, with multiplicity over ratios; the output in
\ref{cls:it:c-ratio} gives the counts for every ratio. The divisor searches use exact
integer arithmetic and a proof-enabled build of GNU \code{factor}, and each list of
factors is checked against its product. The residue masks and the sets of eligible pairs
were reconstructed independently \ref{cls:it:c-review}.
\end{proof}

The use of ratios has the following effect. If one coefficient equals $2c$ and another
$3c$, neither $c$ nor the discriminant need be known: the pair forces a factorization of
$\abs{4q^7-9p^7}$, which depends only on the primes, and one of its divisors would have to
produce an admissible $M$. A single calculation therefore rules out every value of $c$.
The limit $12$ on the reduced multipliers is a choice of finite search size; the interval
inequalities below show that it suffices.

\subsubsection*{Charging the split primes}
Let $\clsXz\le A<x\le U$ and put
\[
  \beta_p=\floor[\big]{\sqrt{4p^7/A}},\qquad \beta=\max_{p\le Y}\beta_p .
\]
Equation \eqref{cls:eq:norm7} implies $b_p\le\beta_p$. Let $k\in\set{1,12}$. Process the
integers $1,\dots,\beta$ in order; at the first unassigned $b$ form a block consisting of
all still unassigned members of $\set{b,2b,\dots,kb}$ that are at most $\beta$. Let $I$
be the set of block minima, and define
\begin{equation}\label{cls:eq:pack}
 L(A,Y,k)=\sum_{b\in I}\ \max_{p\le Y,\ \beta_p\ge b}W_Y(p) .
\end{equation}

\begin{lemma}[Charge bound; \stU]\label{cls:lem:charge}
Let $K$ be as above with $A<x$, and suppose that either $k=1$ and $Y\le6500$, or $k=12$
and $Y\le4200$. Then
\[
   \sum_{\substack{p\le Y\\ p\text{ split}}}W_Y(p)\le L(A,Y,k).
\]
\end{lemma}

\begin{proof}
Every split prime $p\le Y$ has $b_p\le\beta_p\le\beta$, so $b_p$ lies in exactly one
block. Two members of the same block have a ratio whose reduced numerator and denominator
are at most $k$. By \cref{cls:lem:sevenpairs} at most one split prime $p\le Y$ has its
coefficient $b_p$ in any given block; this includes the case of two equal coefficients.
If the minimum of that block is $b$, the prime satisfies $\beta_p\ge b_p\ge b$, so its
charge $W_Y(p)$ is at most the term of \eqref{cls:eq:pack} belonging to the block. Each
term is a separate maximum, so the bound remains valid when the same prime realizes
several maxima, and a block used by no split prime contributes a nonnegative term.
\end{proof}

\subsubsection*{Closing the interval}

\begin{proof}[Proof of \cref{cls:thm:seven}]
Let $K$ be as above. By \ref{cls:it:bound} and \cref{src:thm:dgrh},
$x\le3.9\cdot10^{30}$. Every split prime exceeds $693$, and its square exceeds every $Y$
used below, while no prime $p\le Y$ is ramified. Hence \cref{cls:lem:baseline} gives, for
every window,
\begin{equation}\label{cls:eq:seven}
   \log x>S(Y)-\sum_{\substack{p\le Y\\ p\text{ split}}}W_Y(p).
\end{equation}
Each row of the following table excludes the interval $A<x\le U$. In the first four rows
$A=4Y^7$: if $x>4Y^7$ then every prime up to $Y$ is inert by \cref{cls:cor:witness}(i),
the sum in \eqref{cls:eq:seven} is empty, and $\log x>S(Y)>\log U$ contradicts $x\le U$.
In the remaining rows the sum is bounded by $L(A,Y,k)$ by \cref{cls:lem:charge}, and
$S(Y)-L(A,Y,k)>\log U$ gives the contradiction. The first six rows use the heights $V_0$,
the last eight the heights $V_1$. The entries are outward roundings of the rational
bounds of \cref{cls:rem:exact}: the charge is rounded up, the forced lower bound
$S(Y)-L(A,Y,k)$ down, and $\log U$ up.
\begin{center}
\begin{tabular}{rcllrrr}\toprule
$Y$ & $k$ & $A$ & $U$ & Charge $\le$ & Forced $\log x>$ & $\log U<$\\\midrule
7000 & -- & $4\cdot7000^7$ & $3.9\cdot10^{30}$ & 0 & 76.072 & 70.439\\
5000 & -- & $4\cdot5000^7$ & $4\cdot7000^7$ & 0 & 63.631 & 63.362\\
4700 & -- & $4\cdot4700^7$ & $4\cdot5000^7$ & 0 & 61.571 & 61.007\\
4600 & -- & $4\cdot4600^7$ & $4\cdot4700^7$ & 0 & 60.867 & 60.574\\
6500 & 1 & $4\cdot10^{22}$ & $4\cdot4600^7$ & 12.055 & 61.084 & 60.423\\
4900 & 1 & $10^{22}$ & $4\cdot10^{22}$ & 10.186 & 52.767 & 52.044\\
4200 & 12 & $3.1\cdot10^{20}$ & $3.255\cdot10^{20}$ & 9.223 & 47.295 & 47.232\\
4200 & 12 & $3.255\cdot10^{20}$ & $3.74325\cdot10^{20}$ & 9.001 & 47.518 & 47.372\\
4200 & 12 & $3.74325\cdot10^{20}$ & $4.866225\cdot10^{20}$ & 8.368 & 48.151 & 47.635\\
4200 & 12 & $4.866225\cdot10^{20}$ & $7.2993375\cdot10^{20}$ & 7.322 & 49.197 & 48.040\\
4200 & 12 & $7.2993375\cdot10^{20}$ & $1.31388075\cdot10^{21}$ & 6.011 & 50.508 & 48.628\\
4200 & 12 & $1.31388075\cdot10^{21}$ & $2.6277615\cdot10^{21}$ & 4.600 & 51.919 & 49.321\\
4200 & 12 & $2.6277615\cdot10^{21}$ & $7.8832845\cdot10^{21}$ & 3.425 & 53.094 & 50.420\\
4200 & 12 & $7.8832845\cdot10^{21}$ & $10^{22}$ & 2.120 & 54.399 & 50.657\\\bottomrule
\end{tabular}
\end{center}
The fourteen intervals cover every $x$ with $\clsXz<x\le3.9\cdot10^{30}$. The closest
comparison is the seventh row: at $A=\clsXz$ there are $193$ blocks, and even if every
block uses the most expensive prime available to it the total charge is below $9.223$,
while the inert baseline $S(4200)$ exceeds $56.518$; the difference exceeds the upper
bound $47.232$ for $\log x$. A separate computation at 65-digit working precision
reproduces the fourteen margins, the smallest being $0.0641$ \ref{cls:it:c-explicit}.

Thus no field with $x>\clsXz$ has class group of exponent seven, and \ref{cls:it:census}
gives the $33$ fields. For the orders apply
\cite[Theorem~1.3(6) and Section~6]{Aiken2026Boundary}. Its conductor enumeration starts
from the maximal orders whose class-group exponent divides seven; the fields of class
number one are known, and the result on fields discharges the remaining completeness
condition. This gives the seven nonmaximal orders of \cref{cls:tab:seven}.
\end{proof}

\begin{table}[ht]
\centering
\begin{tabular}{lp{0.70\linewidth}}\toprule
Class group & $\abs{D}$ for the maximal orders\\\midrule
$\Z/7\Z$ & 71, 151, 223, 251, 463, 467, 487, 587, 811, 827, 859, 1163, 1171, 1483, 1523,
1627, 1787, 1987, 2011, 2083, 2179, 2251, 2467, 2707, 3019, 3067, 3187, 3907, 4603, 5107,
5923 \\[3pt]
$(\Z/7\Z)^2$ & 63499, 118843 \\\bottomrule
\end{tabular}

\medskip
\begin{tabular}{rrr}\toprule
Nonmaximal order discriminant & Field discriminant & Conductor\\\midrule
$-284$ & $-71$ &2\\
$-343$ & $-7$ &7\\
$-604$ & $-151$ &2\\
$-892$ & $-223$ &2\\
$-1372$ & $-7$ &14\\
$-1852$ & $-463$ &2\\
$-1948$ & $-487$ &2\\\bottomrule
\end{tabular}
\caption{The $33$ field discriminants and the seven additional order discriminants of
class-group exponent seven. In the lower table the order discriminant is the field
discriminant times the square of the conductor, and the proper class group is cyclic of
order seven.}\label{cls:tab:seven}
\end{table}

\begin{remark}[Verification of the lists]\label{cls:rem:sevenlists}
The $38$ cyclic groups and two groups of rank two in \cref{cls:tab:seven} have together
$364$ classes. A reduced-form computation as in \cref{cls:rem:forms} checks the seventh
power of each class and records generators, together with field discriminant, conductor,
class number and group for each order. A second computation enumerates all $59{,}421$
negative order discriminants of absolute value at most $118{,}843$: of the $135$ orders
whose class number is a nontrivial power of seven, exactly the $40$ recorded here have
exponent seven. Both are in \ref{cls:it:c-sevenlists}.
\end{remark}

The delivered certificates for this subsection are the following.
\begin{enumerate}[resume*=clscert]
\item\label{cls:it:c-explicit} Rational enclosures for the fourteen rows of the table in
the proof of \cref{cls:thm:seven}, their output, and the 65-digit recomputation:\\
\artifact{exponent_seven_and_grh_transfers/checks/explicit_bound.py}\\
\artifact{exponent_seven_and_grh_transfers/evidence/explicit-result.txt}\\
\artifact{exponent_seven_and_grh_transfers/evidence/review_bounds.txt}
\item\label{cls:it:c-equal} The search of \cref{cls:lem:sevenpairs}(i) and its output:\\
\artifact{exponent_seven_and_grh_transfers/checks/equal_norms.py}\\
\artifact{exponent_seven_and_grh_transfers/evidence/equal-result.txt}
\item\label{cls:it:c-ratio} The search of \cref{cls:lem:sevenpairs}(ii), its output, and
the record of the factoring program:\\
\artifact{exponent_seven_and_grh_transfers/checks/ratio_norms.py}\\
\artifact{exponent_seven_and_grh_transfers/evidence/ratio-result.txt}\\
\artifact{exponent_seven_and_grh_transfers/evidence/factor_version.txt}
\item\label{cls:it:c-review} Independent reconstruction of the residue masks and of a
sample of the divisor searches:\\
\artifact{exponent_seven_and_grh_transfers/evidence/review_masks.txt}\\
\artifact{exponent_seven_and_grh_transfers/evidence/crosscheck_output.txt}
\item\label{cls:it:c-sevenlists} The $40$ orders with generators, and the enumeration of
all orders through $118{,}843$:\\
\artifact{exponent_seven_and_grh_transfers/class_groups.csv}\\
\artifact{exponent_seven_and_grh_transfers/evidence/small_range_output.txt}
\end{enumerate}

\subsection{Two general restrictions for even exponents}\label{cls:sec:genus}
\sources{grh_transfers; grh_method_transfers_2026-11-25_1}

The following two restrictions hold for every even exponent and are not used in the
censuses above. For the exponent six they give necessary conditions on a field with
$x>\clsXz$, in terms of the number and size of the prime factors of $x$; those conditions
are contained in \cref{cls:thm:six}, which shows that no such field exists.

\begin{lemma}[Largest prime factor; \stU]\label{cls:lem:pplus}
In the situation of \cref{cls:lem:norm} suppose $j=2r$ is even. Then every prime factor
of $x$ is less than $4p^r$.
\end{lemma}

\begin{proof}
In \eqref{cls:eq:norm} we have $0\le a<2p^r$. The positive integers $2p^r-a$ and
$2p^r+a$ have product $xb^2$, and each is strictly less than $4p^r$. Every prime factor
of $x$ divides at least one of them.
\end{proof}

For the second restriction let $D<-4$, write $A=\Cl(K)$ and $t=\omega(x)$, and call a
squarefree product of the prime ideals above the rational primes dividing $D$ a
\emph{ramified ideal}. The genus field $G$ of $K$ is the compositum of the quadratic
fields attached to the $t$ prime discriminants dividing $D$. It is multiquadratic of
degree $2^t$ over $\Q$, unramified over $K$ at every finite prime, and
$\Gal(G/K)\simeq A/A^2$ by the principal genus theorem \cite[\S6]{Cox2013Primes}.

\begin{theorem}[Intermediate genus-field bound; \stG]\label{cls:thm:genus}
Suppose the exponent of $A$ divides $2m$, where $m$ is a positive odd integer. Let
$H\subseteq A[2]$ have dimension $k$ over $\F_2$, and assume that each class in $H$ has a
ramified ideal representative of norm at most $R$. Then there is an abelian field
$F/\Q$ containing $K$, of degree $n=2^{t-k}$ and absolute discriminant $x^{n/2}$, such
that every unramified rational prime $p$ splitting completely in $F$ satisfies
\begin{equation}\label{cls:eq:genusnorm}
                         x\le4Rp^m .
\end{equation}
Consequently
\begin{equation}\label{cls:eq:genusuniv}
 x\le4R\bigl(2n\log x+\tfrac52n+5\bigr)^{2m}.
\end{equation}
The field $F$ and the inequality \eqref{cls:eq:genusnorm} are unconditional;
\eqref{cls:eq:genusuniv} uses \cref{src:thm:dgrh}. Any sharper valid upper bound for the
least completely split prime of $F$ may replace the one used in
\eqref{cls:eq:genusuniv}.
\end{theorem}

\begin{proof}
The map $A[2]\to A/A^2$ is an isomorphism. Indeed $A^2$ has exponent dividing the odd
integer $m$, so it meets $A[2]$ trivially; and for every $a\in A$ the element $a^m$ lies
in $A[2]$ while $a^{1-m}$ is a square. Regard $H$ as a subgroup of $\Gal(G/K)$ and put
$F=G^H$. Then $K\subseteq F$, $[F:\Q]=n$, and $F/K$ is unramified, so the discriminant
tower formula gives $\abs{\Delta_F}=\abs{D}^{[F:K]}=x^{n/2}$. Moreover $F/\Q$ is abelian
because $G/\Q$ is.

Let $p$ be as stated. Since $F\supseteq K$, the prime $p$ splits in $K$ and is prime to
$D$. Choose a prime $\clsp$ of $K$ above $p$. Its Artin symbol in $\Gal(G/K)$ lies in $H$,
so $[\clsp]=ha^2$ for some $h\in H$ and $a\in A$, and raising to the $m$th power gives
$[\clsp]^m=h$. Choose a ramified ideal $J$ in the class $h$ with $\Nm(J)\le R$. Because
$h^2=1$, the integral ideal $\clsp^mJ$ is principal, say $(\alpha)$. Its valuations at
$\clsp$ and $\bar\clsp$ are $m$ and zero, since $J$ is supported at primes above $D$.
Hence $\alpha\notin\Q$. Writing $\alpha=(a+b\sqrt D)/2$ with integers $a,b$ and $b\ne0$,
norms give
\[
 p^m\Nm(J)=\Nm_{K/\Q}(\alpha)=\frac{a^2+xb^2}{4}\ge\frac x4 ,
\]
which proves \eqref{cls:eq:genusnorm}. The cost in the norm is $R$, not $R^m$, because the
ramified ideal is inserted after the $m$th power is taken.

The Dedekind zeta function of the abelian field $F$ is a product of Dirichlet
$L$-functions, so \cref{src:thm:dgrh} gives its Riemann hypothesis. Theorem~5.1 of
\cite{Bach1996Explicit} then provides an unramified prime with trivial Frobenius in
$\Gal(F/\Q)$ and
\[
 p\le\bigl(4\log\abs{\Delta_F}+\tfrac52n+5\bigr)^2 .
\]
Substitution into \eqref{cls:eq:genusnorm} proves \eqref{cls:eq:genusuniv}.
\end{proof}

One may optimize the theorem by replacing $R$ with
\[
 R(H)=\max_{h\in H}\min\set{\Nm(J):J\text{ ramified},\ [J]=h}
\]
and varying $H$. The next
lemma gives an explicit choice that avoids any search over subgroups. Let
$q_1<\dots<q_t$ be the rational primes dividing $D$, with ramified prime ideals
$\clsq_i$, and write $P_j=q_1\cdots q_j$, $P_0=1$.

\begin{lemma}[Independence of ramified classes; \stU]\label{cls:lem:ramified}
Under the exponent hypothesis of \cref{cls:thm:genus}, any selection of at most $t-2$
ramified prime classes is independent in $A[2]$ and in $A/A^2$. Thus for $0\le k\le t-2$
the subgroup generated by $[\clsq_1],\dots,[\clsq_k]$ has dimension $k$ and has ramified
representatives of norm at most $P_k$.
\end{lemma}

\begin{proof}
Each ramified class has order dividing two. Suppose a nonempty squarefree product $J$ of
ramified prime ideals is principal, say $J=(\alpha)$. Conjugation fixes $J$, so
$\bar\alpha/\alpha$ is a unit of $K$, and since $D<-4$ the units are $\pm1$. If
$\bar\alpha=\alpha$, then $\alpha$ is rational and its valuations at all ramified primes
are even, contrary to the definition of $J$. Otherwise $\alpha$ is a rational multiple of
$\sqrt D$. Multiplication by a rational number changes every ramified valuation by an even
integer. At each odd prime dividing $D$ the valuation of $\sqrt D$ is odd. At the prime
above two its parity is that of $v_2(D)$, and a fundamental even discriminant has
$v_2(D)=2$ or $3$. Therefore a nonempty relation is supported on all $t$ ramified primes
when $D$ is odd or $v_2(D)=3$, and on all $t-1$ odd ramified primes when $v_2(D)=2$. No
selection of at most $t-2$ ramified primes supports a relation. Finally, a class of order
dividing two that lies in $A^2$ is trivial, since $A^2$ has odd exponent; this proves
independence in the genus quotient as well. Every class of the stated subgroup has a
representative formed from a subset of the first $k$ prime ideals, of norm at most $P_k$.
\end{proof}

\subsection{Open questions}\label{cls:sec:open}
\sources{exponent_seven_and_grh_transfers; grh-exponent-six-and-transfers}

\begin{enumerate}
\item In the proof of \cref{cls:thm:seven} the reduced multipliers are limited to $12$.
What is the smallest maximal reduced multiplier that suffices after optimizing the test
function and the partition into blocks?
\item Can the pair exclusions of \cref{cls:lem:sixpairs,cls:lem:sevenpairs} be verified by
a shorter, formally checked arithmetic certificate? Neither census has been formalized;
both rest on \cref{src:thm:dgrh}, on the published enumeration \ref{cls:it:census} and on
the finite computations cited above.
\item For exponents above eight the same norm equations can be tested, but a new initial
census range, new prime weights and a new compatibility proof for pairs of equations would
be needed. For which further exponents can the interval between an enumerated range and
the bound of type \ref{cls:it:bound} be closed in this way?
\end{enumerate}
